\documentclass[11pt,reqno]{amsart}

\usepackage[a4paper,margin=1.15in]{geometry}
\usepackage[T1]{fontenc}
\usepackage{lmodern}
\usepackage{amsmath,amssymb,mathtools}
\usepackage{microtype}
\usepackage{enumitem}
\usepackage[hidelinks]{hyperref}
\raggedbottom
\hypersetup{
  pdftitle={Stabilizer/GNS Support Realization and Finite Central Fibers},
  pdfauthor={}
}

\numberwithin{equation}{section}

\newtheorem{theorem}{Theorem}[section]
\newtheorem{proposition}[theorem]{Proposition}
\newtheorem{lemma}[theorem]{Lemma}
\newtheorem{corollary}[theorem]{Corollary}
\theoremstyle{definition}
\newtheorem{definition}[theorem]{Definition}
\newtheorem{assumption}[theorem]{Assumption}
\theoremstyle{remark}
\newtheorem{remark}[theorem]{Remark}

\newcommand{\cA}{\mathcal A}
\newcommand{\cB}{\mathcal B}
\newcommand{\cH}{\mathcal H}
\newcommand{\cM}{\mathcal M}
\newcommand{\cR}{\mathcal R}
\newcommand{\cC}{\mathcal C}
\newcommand{\cZ}{\mathcal Z}
\newcommand{\id}{\mathbf 1}
\newcommand{\Tr}{\operatorname{Tr}}
\newcommand{\supp}{\operatorname{supp}}
\newcommand{\fin}{\mathrm{fin}}
\newcommand{\sot}{\mathrm{s}}

\title[Stabilizer/GNS Support Realization]{Stabilizer/GNS Support Realization and Finite Central Fibers\\
\large A checked bridge from finite stabilizer averages and specified GNS-orbit hypotheses}
\author{}
\date{August 3, 2026}

\begin{document}

\begin{abstract}
This note records the stabilizer/GNS layer behind the finite-fiber theorem in
the current stabilizer-region manuscript.  It machine-checks four load-bearing
implications that the earlier abstract corner theorem left open: a finite
self-adjoint group average is a projection; nested finite stabilizer subgroups
give decreasing local supports; every such decreasing sequence has a canonical
strong limit; and the specified dense regional/complementary product orbit
makes the compressed vector functional projection-faithful and the support
full.  Combined with the checked full-tracial-corner theorem of
Investigation~I001, this proves that every nonzero central summand contains a
nonzero finite projection, and supplies the canonical witness in every
explicitly presented finite family of central fibers.

The note is deliberately fail-closed about the remaining analytic layer.  The
Lean theorem constructs the strong limit after receiving the finite group
averages, but it still receives the represented stabilizer action, the
stabilized GNS vector, and the already-extended corner-traciality statement.
It does not yet construct the tensor-product Pauli realization, the
quasi-local stabilizer state or its reduced density matrices, derive the
Kaplansky extension of local traciality, construct the GNS product orbit, or
formalize direct-integral/type-III vocabulary.  Those are
reported as open formalization obligations rather than hidden axioms.
\end{abstract}

\maketitle

\section{Target and checked scope}
\label{sec:target}

Let $\cM\subseteq\cB(\cH)$ be a concrete von Neumann algebra, let
$\xi\in\cH$, and let $p_n\in\cM$ be local support projections.  The target
is the central-summand conclusion of the manuscript's theorem
\texttt{thm:finite\_projections\_exist}:
\begin{equation}
  0\neq z\in\cZ(\cM)
  \quad\Longrightarrow\quad
  \exists e\in\cM:\ e\neq0,\ e\leq z,\ e\text{ is finite}.
  \label{eq:central_summand_goal}
\end{equation}
The final phrase is the Murray--von Neumann finiteness of a projection in the
ambient concrete algebra.  Since $e\leq z$, it is the concrete
central-summand statement that $z\cM$ contains a finite projection.  This is
stronger and more precise than merely presenting a finite list of fibers.

The formal development in
\texttt{formalization\_finite\_fibers/FiniteFiber/} now proves the following
conditional theorem without project-local axioms or placeholders.

\begin{theorem}[Checked nested-stabilizer-sequence-to-central-summand theorem]
\label{thm:checked_stabilizer_to_central_summand}
Assume the following data.
\begin{enumerate}[label=\textup{(H\arabic*)},leftmargin=3.1em]
  \item A group $K$ acts on $\cH$ by bounded operators $\rho(k)$, every
  $\rho(k)$ is self-adjoint, and $S_n\leq K$ is a finite increasing sequence
  of subgroups.  Define
  \begin{equation}
    p_n=\frac{1}{\lvert S_n\rvert}\sum_{s\in S_n}\rho(s).
    \label{eq:stabilizer_group_average}
  \end{equation}
  \item The represented finite stabilizers lie in $\cM$ and fix $\xi$:
  \begin{equation}
    \rho(s)\in\cM,
    \qquad
    \rho(s)\xi=\xi,
    \qquad s\in S_n.
    \label{eq:represented_stabilizers_fix_xi}
  \end{equation}
  The formal construction defines the strong limit $q$ from the $p_n$.
  \item There are regional and complementary operator sets
  $\cR,\cC\subseteq\cB(\cH)$ such that $\cR\subseteq\cM$,
  region and complement commute, every element of $\cM$ commutes with
  $\cC$, and
  \begin{equation}
    \overline{\operatorname{span}\{ab\xi:a\in\cR,\ b\in\cC\}}=\cH.
    \label{eq:product_orbit_dense}
  \end{equation}
  \item The compressed vector functional
  \begin{equation}
    \tau(x):=\langle\xi,x\xi\rangle
    \label{eq:compressed_vector_functional}
  \end{equation}
  is tracial on the support corner: for $a,b\in q\cM q$,
  \begin{equation}
    \tau(ab)=\tau(ba).
    \label{eq:corner_traciality}
  \end{equation}
\end{enumerate}
Then $q\in\cM$ is a full projection, and its tracial vector functional is
faithful on projections below $q$ (the exact property used for the
Murray--von Neumann finiteness argument); \eqref{eq:central_summand_goal}
holds.  More explicitly, the witness for a nonzero central projection $z$ is
\begin{equation}
  e=zq.
  \label{eq:canonical_central_witness}
\end{equation}
For any supplied finite pairwise-orthogonal central partition
$(z_i)_{i\in I}$ of $\id$, each $z_iq$ is therefore a nonzero finite
projection below $z_i$.
\end{theorem}

The exact Lean entry point for all clauses of
Theorem~\ref{thm:checked_stabilizer_to_central_summand} is the following
declaration.  Its fields \texttt{q\_projection},
\texttt{q\_fixes\_xi}, \texttt{full}, \texttt{projection\_faithful}, and
\texttt{central\_witness} match the displayed conclusions; \texttt{semifinite}
is their existential central-summand corollary.  The line segments concatenate
within the identifier:
\begin{quote}\small
\texttt{stabilizerAverageSequence\_}\\
\texttt{canonicalSupportCornerConclusion\_}\\
\texttt{of\_gnsOrbit}
\end{quote}
The canonical semifiniteness and central-witness specializations have the
separate entry points
\begin{quote}\small
\texttt{stabilizerAverageSequence\_canonicalSemifinite\_of\_gnsOrbit}\\
\texttt{stabilizerAverageSequence\_canonicalCentralWitness\_of\_gnsOrbit}.
\end{quote}

\section{Strong-limit support}
\label{sec:strong_support}

The formal predicate \texttt{StronglyTendsTo} means pointwise norm
convergence on $\cH$.  The development no longer receives it as a hypothesis.
For $U_n=(\operatorname{ran}p_n)^\perp$, it defines
\begin{equation}
  q=\id-\operatorname{Proj}_{\overline{\bigcup_n U_n}}.
  \label{eq:formal_strong_limit}
\end{equation}
The ranges are decreasing because $S_n\leq S_m$ implies that the average over
$S_m$ is a subprojection of the average over $S_n$.  Thus the $U_n$ are
increasing, and mathlib's theorem on monotone orthogonal projections gives
pointwise norm convergence to \eqref{eq:formal_strong_limit}.

\begin{proposition}[Constructed strong support]
\label{prop:strong_support_consequences}
Under (H1)--(H2), each $p_n$ is a projection in $\cM$, the sequence is
decreasing, $p_n\xi=\xi$, and the constructed $q$ is a projection in $\cM$
with $q\xi=\xi$.
\end{proposition}

\begin{proof}
For a finite group average, left and right multiplication by a represented
group element permutes the finite sum.  Hence the average is idempotent; the
self-adjointness of every summand makes it a projection.  If $S_n\leq S_m$,
each summand of $p_n$ fixes $p_m$ on both sides, so
\begin{equation}
  p_mp_n=p_np_m=p_m.
  \label{eq:stage_fixes_limit}
\end{equation}
Every element of $S_n$ fixes $\xi$, so averaging gives $p_n\xi=\xi$.  The
orthogonal-complement construction in \eqref{eq:formal_strong_limit} gives
$p_n\to q$ strongly.  The checked strong-limit consequences then prove
$q\in\cM$, $q^2=q=q^*$, and $q\xi=\xi$.
\end{proof}

The construction is compatible with the standard operator-algebraic result
that a decreasing family of projections has a strong limit; compare
Blackadar, I.5.1.3, PDF p.~29.  The strong-closure and corner facts used in
the manuscript are also located in Blackadar, I.9.1.2--I.9.1.4, PDF p.~58.
The Lean proof does not appeal to an unimplemented closure axiom: it builds
the limit through orthogonal projections and derives membership directly from
the concrete double-commutant definition.

\section{Projection-faithfulness from a specified GNS product orbit}
\label{sec:faithfulness}

The key correction to the first formalization layer is that
projection-faithfulness is not now a field of the theorem statement.  It is
derived from (H3)--(H4).
Let $r\leq q$ be a projection satisfying $\tau(r)=0$.  Since $r=r^*=r^2$,
\begin{equation}
  0=\tau(r)=\langle r\xi,r\xi\rangle,
  \qquad\text{so}\qquad r\xi=0.
  \label{eq:null_projection_annihilates_xi}
\end{equation}
For $y\in q\cM q$, all three factors below remain in the corner, and
traciality gives
\begin{equation}
  \begin{aligned}
    \lVert ry\xi\rVert^{2}
    &=\tau(y^*ry) \\
    &=\tau(yy^*r) \\
    &=0.
  \end{aligned}
  \label{eq:corner_nullity}
\end{equation}
For $a\in\cR$, choose $y=qaq$.  Equations
\eqref{eq:corner_nullity} and $q\xi=\xi$ show that $ra\xi=0$.  If
$b\in\cC$, locality and the commutation of $\cM$ with $\cC$ yield
\begin{equation}
  rab\xi=bra\xi=0.
  \label{eq:product_orbit_annihilation}
\end{equation}
Density in \eqref{eq:product_orbit_dense} and continuity of $r$ now imply
$r=0$.  Thus the vector functional is faithful on all support
subprojections, which is exactly the projection-faithfulness needed by the
Murray--von Neumann finiteness argument.  The stronger positive-element
faithfulness of Proposition~\texttt{prop:faithful} in the manuscript is not
claimed as a Lean theorem here.

The corresponding checked declaration is the concatenation
\begin{quote}\small
\texttt{VectorCornerData.}\\
\hspace*{1em}\texttt{projection\_eq\_zero\_}\\
\hspace*{1em}\texttt{of\_vectorFunctional\_eq\_zero}.
\end{quote}
The proof exposes the two commutations that cannot be omitted: regional with
complement, and $\cM$ with complement.

\section{Fullness and the central witness}
\label{sec:fullness}

Let $z\in\cZ(\cM)$ and suppose $zq=0$.  Since $q\xi=\xi$,
\begin{equation}
  z\xi=zq\xi=0.
  \label{eq:central_nullity}
\end{equation}
Centrality commutes $z$ with regional operators; (H3) commutes $z$ with the
complement.  Hence $zab\xi=abz\xi=0$ for every product-orbit vector.
Equation~\eqref{eq:product_orbit_dense} forces $z=0$.  Therefore $q$ is full
in the operational form
\begin{equation}
  0\neq z\in\cZ(\cM)
  \quad\Longrightarrow\quad
  zq\neq0.
  \label{eq:checked_fullness}
\end{equation}

By the checked I001 theorem, a faithful tracial support corner makes every
subprojection of $q$ finite.  For nonzero central $z$, the product $zq$ is a
projection, is below $q$ and $z$, and is nonzero by
\eqref{eq:checked_fullness}.  It is thus the claimed finite witness in
\eqref{eq:canonical_central_witness}.  This is the algebraic content of the
central-carrier argument; compare Blackadar, III.1.1.5, PDF p.~245, the
definition of finite projection, III.1.3.1, PDF p.~250, and the
semifinite/purely-infinite distinction, III.1.4.2, PDF p.~253.

\section{Exact boundary of the checked theorem}
\label{sec:boundary}

The theorem above is a faithful formalization of the finite-group/support-
corner portion of the live proof, but it is not an end-to-end construction
from a raw tensor-product Pauli presentation.  Five boundaries matter.

\begin{enumerate}[label=\textup{(B\arabic*)},leftmargin=3.1em]
  \item \textbf{Pauli and density-matrix realization.}  The checked finite
  group layer proves the projection and monotonicity properties of
  \eqref{eq:stabilizer_group_average} from an explicitly supplied
  self-adjoint operator action and a nested finite-subgroup sequence.  Its
  separate finite-dimensional trace lemma assumes, exactly,
  $\Tr(\rho(s))=0$ for every $s\neq1$ and $\dim\cH_F\neq0$; under those
  hypotheses it proves $\Tr(P_F)=\dim\cH_F/\lvert S_F\rvert$ and
  $P_F\neq0$.  That lemma is not used to produce a nonzero global GNS
  support in the central-summand theorem.  The package does not yet construct
  the action, the tracelessness condition, or the qubit-dimension identity
  from tensor-product Pauli operators, and it does not prove partial-trace
  consistency of the reduced density matrices.
  \item \textbf{Stabilizer-state realization.}  The condition
  \eqref{eq:represented_stabilizers_fix_xi} is explicit.  The Lean package
  does not yet construct the quasi-local stabilizer state and prove that its
  GNS vector has this property.
  \item \textbf{GNS-orbit realization.}  The structure
  \texttt{GNSOrbitConditions} receives the regional and complementary operator
  sets, their membership and commutation properties, and the density statement
  \eqref{eq:product_orbit_dense} as fields.  It packages their exact
  consequence for the support corner, but does not construct the GNS
  representation or prove density from the quasi-local algebra.  This is a
  separate formalization task; in the installed mathlib version the relevant
  GNS source still records the cyclic-vector construction as unfinished.
  \item \textbf{Local trace extension.}  The manuscript obtains the
  corner-traciality hypothesis by
  finite-dimensional matrix-trace cyclicity, bounded strong continuity of
  multiplication, and Kaplansky density.  Blackadar states that multiplication
  is separately continuous in all of the listed operator topologies and jointly
  continuous on bounded sets for the strong and strong-$*$ topologies in
  I.3.2.1, PDF p.~23; Kaplansky density is I.9.1.3, PDF p.~58.
  Mathlib~v4.28 has no ready-made theorem connecting those ingredients to this
  corner-traciality statement, so it remains an explicit analytic input.
  \item \textbf{Fiber/type language.}  The installed mathlib version has no
  direct-integral or type-III-factor API.  The checked result is therefore
  the stronger central-projection statement \eqref{eq:central_summand_goal}
  and its finite-family specialization, rather than an unimplemented
  ``almost every direct-integral factor'' predicate.
\end{enumerate}

These boundaries are not defects in the theorem statement: they prevent the
formalization from silently converting an analytic or measure-theoretic step
into a project-local axiom.  In particular, the theorem does not state an
unchecked \texttt{not\_typeIII} conclusion.

\section{Lean artifact and verification map}
\label{sec:verification}

\begin{center}\footnotesize
\begin{tabular}{p{.34\linewidth}p{.58\linewidth}}
\hline
Lean module & Checked responsibility \\
\hline
\texttt{FiniteFiber/Basic.lean} & projections, central slices, Murray--von Neumann finiteness \\
\texttt{FiniteFiber/TracialCorner.lean} & a projection-faithful tracial corner makes support subprojections finite \\
\texttt{FiniteFiber/FiniteFibers.lean} & full support gives finite witnesses below all nonzero central projections \\
\texttt{FiniteFiber/}\par\texttt{FiniteStabilizer.lean} & finite group averages are projections and respect nested subgroup constraints \\
\texttt{FiniteFiber/}\par\texttt{FiniteStabilizerTrace.lean} & conditional finite-dimensional trace/rank formula and nonzero local support \\
\texttt{FiniteFiber/}\par\texttt{StabilizerSequence.lean} & an increasing finite-subgroup sequence gives decreasing supports and fixed vectors \\
\texttt{FiniteFiber/StrongSupport.lean} & constructs the strong limit and proves its membership, projection, and fixed-vector properties \\
\texttt{FiniteFiber/VectorCorner.lean} & projection-faithfulness and fullness from the dense product orbit \\
\texttt{FiniteFiber/GNSOrbit.lean} & bundled I002 conclusion and finite-family witness \\
\texttt{FiniteFiber/}\par\texttt{StabilizerFiniteFiber.lean} & central-fiber conclusion from nested subgroup data \\
\hline
\end{tabular}
\end{center}

The expected verification commands are
\begin{verbatim}
lake build FiniteFiber.StrongSupport
lake build FiniteFiber.FiniteStabilizer
lake build FiniteFiber.FiniteStabilizerTrace
lake build FiniteFiber.StabilizerSequence
lake build FiniteFiber.VectorCorner
lake build FiniteFiber.GNSOrbit
lake build FiniteFiber.StabilizerFiniteFiber
lake build
lake env lean tests/AxiomAudit.lean
\end{verbatim}
The axiom audit must report only Lean's three standard logical axioms for all
public theorems, including the bundled I002 theorem and finite-family witness.
\begin{center}
\texttt{propext}, \texttt{Classical.choice}, \texttt{Quot.sound}
\end{center}

\section{Source locator policy}
\label{sec:sources}

The staged PDF used for the operator-algebra references is
\texttt{sources/blackadar.pdf}.  Citations in this note use printed
section/theorem numbers together with PDF page numbers.  The mathematical
derivations in Sections~\ref{sec:strong_support}--\ref{sec:fullness} are
proved directly in Lean; the citations identify the standard analytic and
classification facts delimiting the formal interface.

\end{document}
